\begin{name}
	{\tenchude}
	{TOÁN 10}
	{LỚP TOÁN THẦY PHÁT}
	{Thất bại đơn giản chỉ là cơ hội để ta bắt đầu lại mọi thứ một cách thông minh hơn.}
\end{name}
\Opensolutionfile{ansbook}[ans/ansbookDe2]
\TN
\Opensolutionfile{ans}[ans/ansDe2-TN1]
\begin{ex}%[0-HK1-CT-6-HungVuong-HCM-2324]%[VN-MT-6, Nguyễn Hữu Đức]%[0D6N3-3]
Một nhóm có $10$ học sinh tham gia một kỳ thi. Số điểm thi của $10$ học sinh đó (thang điểm $10$) được sắp xếp từ thấp đến cao như sau:
\begin{center}
$\begin{array}{cccccccccc}
$0$&$1$&$2$&$4$&$4$&$5$&$7$&$8$&$8$&$9$.
\end{array}$
\end{center}
Khi đó, trung vị của mẫu số liệu trên là
\choice
{$5$}
{$4$}
{\True $4{,}5$}
{$5{,}5$}
\loigiai{
Sắp xếp mẫu số liệu theo thứ tự không giảm ta được
\begin{center}
$\begin{array}{cccccccccc}
$0$&$1$&$2$&$4$&$4$&$5$&$7$&$8$&$8$&$9$.
\end{array}$
\end{center}
Số liệu có $10$ phần tử nên trung vị là $M_e=\dfrac{4+5}{2}=4{,}5$.
}
\end{ex}

\begin{ex}%[HKI, Phan Bội Châu- Bình Thuận, 2023]%[Hiếu Phan]%[0D6N4-2]
Cho mẫu số liệu: $4;11;17;18;22;24;30;33;72;87$. Khoảng biến thiên của mẫu số liệu thống kê trên là
\choice
{\True $83$}
{$33$}
{$89$}
{$0$}
\loigiai{ Khoảng biến thiên của mẫu số liệu là $87-4=83$.
}
\end{ex}

\begin{ex}%[10-11-12EX-GK2-2425]%[Trần Văn Hùng]%[0D8H1-3]
\immini[thm]{
Từ thành phố $A$ đến thành phố $B$ có $3$ con đường, từ thành phố $A$ đến thành phố $C$ có $2$ con đường, từ thành phố $B$ đến thành phố $D$ có $2$ con đường, từ thành phố $C$ đến thành phố $D$ có $3$ con đường. Biết rằng không có con đường nào nối trực tiếp từ thành phố $B$ đến thành phố $C$ và không có con đường nào nối trực tiếp từ thành phố $A$ đến thành phố $D$. Hỏi có bao nhiêu cách để đi từ thành phố $A$ đến thành phố $D$?
\choice
{$6$}
{\True $12$}
{$18$}
{$36$}
}
{
\begin{tikzpicture}[scale=1, font=\footnotesize, line join=round, line cap=round, >=stealth]
\coordinate (A) at (0,0);
\coordinate (B) at (2,2);
\coordinate (C) at (2,-2);
\coordinate (D) at (5,0);
\draw (A)--(B)--(D)--(C)--(A)
(1,1) node[above]{$3$} (4,1)node[scale=1]{$2$}
(1,-1) node[above]{$2$} (3.5,-1) node[above]{$3$}
;
\foreach \i/\g in {A/180,B/90,C/-90,D/0}{\draw[fill=black](\i) circle (1.5pt) ($(\i)+(\g:3mm)$) node[scale=1]{$\i$};}
\end{tikzpicture}
}
\loigiai{
Để đi từ $A$ đến $D$ có hai cách đi là đi qua $B$ hoặc đi qua $C$.
\begin{itemize}
\item Số cách đi từ $A$ đến $D$ qua $B$ là $3\cdot 2=6$ cách.
\item Số cách đi từ $A$ đến $D$ qua $C$ là $2\cdot 3=6$ cách.
\end{itemize}
Vậy có $6+6=12$ cách đi từ thành phố $A$ đến thành phố $D$.
}
\end{ex}

\begin{ex}%[0D8H2-3]
Có bao nhiêu số tự nhiên có bảy chữ số khác nhau từng đôi một, trong đó chữ số $ 2$ đứng liền giữa hai chữ số $ 1$ và $3$ ?
\choice
{$3204$ số}
{$249$ số}
{$2942$ số}
{\True $7440$ số}
\loigiai{
Vì chữ số $ 2$ đứng liền giữa hai chữ số $ 1$ và$ 3$ nên số cần lập có bộ ba số $ 123$ hoặc $ 321$.
\begin{itemize}
\item \textbf{Trường hợp 1:} Số cần lập có bộ ba số $ 123$.\\
Nếu bộ ba số $ 123$ đứng đầu thì số có dạng $\overline{123abcd}$.\\
Có $ \mathrm{A}_7^4=840$ cách chọn bốn số $ a$, $ b$, $ c$, $ d$ nên có $ \mathrm{A}_7^4=840$ số.\\
Nếu bộ ba số $ 123$ không đứng đầu thì số có $ 4$ vị trí đặt bộ ba số $ 123$.\\
Có $ 6$ cách chọn số đứng đầu và có $ \mathrm{A}_6^3=120$ cách chọn ba số $ b$, $ c$, $ d$.\\
Theo quy tắc nhân có $ 6\cdot4\cdot\mathrm{A}_6^3=2880$ số\\
Theo quy tắc cộng có $ 840+2880=3720$ số.
\item \textbf{Trường hợp 2:} Số cần lập có bộ ba số $ 321$.
\end{itemize}
Do vai trò của bộ ba số $ 123$ và $321$ như nhau nên có $ 2\left(840+2880\right)=7440$}
\end{ex}

\begin{ex}%[0D8H2-4]
Một tổ công nhân có $ 12$ người. Cần chọn $ 3$ người, một người làm tổ trưởng, một tổ phó và một thành viên. Hỏi có bao nhiêu cách chọn?
\choice
{$ 220$}
{$ 12!$}
{\True $ 1320$}
{$ 1230$}
\loigiai{
Số cách chọn $ 3$ người, một người làm tổ trưởng, một tổ phó và một thành viên là
\[\mathrm{C}_{12}^1\cdot\mathrm{C}_{11}^1\cdot\mathrm{C}_{10}^1=1320.\]
}
\end{ex}

\begin{ex}%[0D8H2-6]%[Dự án đề kiểm tra Toán khối 10 HK2-NH23-24-Đợt 15-Dương Công Tạo]%[THPT Đặng Huy Trứ - Huế]
Từ $18$ điểm phân biệt trong mặt phẳng và không có $3$ điểm nào thẳng hàng, có thể vẽ được bao nhiêu tam giác có $3$ đỉnh là $3$ trong $18$ điểm đã cho?
\choice
{$6$}
{\True $\mathrm{C}_{18}^3$}
{$\mathrm{A}_{18}^3$}
{$\dfrac{18!}{3}$}
\loigiai{
Số tam giác có $3$ đỉnh là $3$ trong $18$ điểm đã cho là tổ hợp chập $3$ của $18$.\\
Vậy ta có $\mathrm{C}_{18}^3$ (tam giác).
}
\end{ex}

\begin{ex}%[0D8H3-2]  %[Dự án đề kiểm tra Toán 10 HKII NH23-24- Thầy Hải Toán]%[THPT Đống Đa Hà Nội]
Biểu thức $81x^4-108x^3+54x^2-12x+1$ là khai triển Niuton của nhị thức nào sau đây?
\choice
{\True $(3x-1)^4$}
{$(x+3)^4$}
{$(x-3)^4$}
{$(3x+1)^4$}
\loigiai{
Ta có \allowdisplaybreaks
\begin{eqnarray*}
(3x-1)^4&=&\mathrm{C}_{4}^{0}\cdot\left(3x\right)^4+\mathrm{C}_{4}^{1}\cdot\left(3x\right)^3\cdot(-1)^1+\mathrm{C}_{4}^{2}\cdot\left(3x\right)^2\cdot(-1)^2+\mathrm{C}_{4}^{3}\cdot\left(3x\right)\cdot(-1)^3+\mathrm{C}_{4}^{4}\cdot(-1)^4\\
&=& 81x^4-108x^3+54x^2-12x+1.
\end{eqnarray*}
}
\end{ex}

\begin{ex}%[De-chuan-hoa-so-13]%[Huỳnh Quy]%[0H9H1-1]
Sự chuyển động của một tàu thủy được thể hiện trên một mặt phẳng tọa độ như sau: Tàu khởi hành từ vị trí $A(1;2)$ chuyển động thẳng đều với vận tốc (tính theo giờ) được biểu thị bởi véc-tơ $\overrightarrow{v}(3;1)$. Tại thời điểm sau khi khởi hành $2$ giờ thì vị trí của tàu là điểm $B$ (trên mặt phẳng tọa độ) có tọa độ là
\choice
{\True $(7;4)$}
{$(4;3)$}
{$(-5;0)$}
{$(5;5)$}
\loigiai{
Sau $2$ giờ vị trí của tàu tại điểm $B(7;4)$.
}
\end{ex}

\begin{ex}%[0H9H3-2]
Đường trung trực của đoạn $AB$ với $A\left( 4;-1 \right)$ và $B\left( 1;-4 \right)$ có phương trình là
\choice
{$x+y=1$}
{\True $x+y=0$}
{$y-x=0$}
{$x-y=1$}
\loigiai{
Gọi $I$ là trung điểm của $AB$ và $d$ là trung trực đoạn $AB$.\\
Từ $A\left( 4;-1 \right),\,B\left( 1;-4 \right)\to I\left( \dfrac{5}{2};-\dfrac{5}{2} \right)\in d$.\\
Do $d\bot AB\to {{{\vec n}}_d}=\vec{AB}=\left( -3;-3 \right)=-3\left( 1;1 \right)$.\\
Vậy phương trình đường trung trực có dạng\\ \[a\left(x-x_0\right)+b\left(y-y_0\right)=0 \Leftrightarrow 1\left(x-\dfrac{5}{2}\right)+1\left(y+\dfrac{5}{2}\right)=0\Leftrightarrow x+y=0.\]
}
\end{ex}

\begin{ex}%[0H9H4-2]%[Dự án đề kiểm tra Toán khối 10 HKI NH23-24-Dot 15- Hồ Văn Trung]%[THPT Lê Qúy Đôn- TPHCM]
Trong mặt phẳng $Oxy$, phương trình đường tròn có tâm $I\left(3;1\right)$ và đi qua điểm $M\left(2;-1\right)$ là
\choice
{$\left(x+3\right)^2+\left(y+1\right)^2=5$}
{$\left(x-3\right)^2+\left(y-1\right)^2=\sqrt{5}$}
{$\left(x+3\right)^2+\left(y+1\right)^2=\sqrt{5}$}
{\True $\left(x-3\right)^2+\left(y-1\right)^2=5$}
\loigiai{Ta có $IM=\sqrt{(2-3)^2+(-1-1)^2}=\sqrt{5}$.\\
Đường tròn tâm $I(3;1)$ và có $R=IM=\sqrt{5}$ là
$(C) \colon (x-3)^2+(y-1)^2=5$.
}
\end{ex}

\begin{ex}%[0H9H5-2]
Trong mặt phẳng $Oxy$, phương trình chính tắc của Elip có độ dài  trục lớn  bằng $6$ và tỉ số giữa tiêu cự với độ dài trục lớn bằng $\dfrac{1}{3}$ là
\choice
{$\dfrac{x^2}{9}+\dfrac{y^2}{3}=1$}
{$\dfrac{x^2}{9}+\dfrac{y^2}{5}=1$}
{\True $\dfrac{x^2}{9}+\dfrac{y^2}{8}=1$}
{$\dfrac{x^2}{6}+\dfrac{y^2}{5}=1$}
\loigiai{ Độ dài trục lớn $2a=6\Rightarrow a=3$.\\
Tỉ số giữa tiêu cự và độ dài trục lớn  $\dfrac{2c}{2a}=\dfrac{1}{3}\Rightarrow c =\dfrac{1}{3} \cdot a=\dfrac{1}{3}\cdot 3=1$.\\
Ta có $c^2=a^2-b^2\Leftrightarrow 1^2=3^2-b^2\Leftrightarrow b^2=3^2-1^2=8$. \\
Phương trình chính tắc của Elip $\dfrac{x^2}{9}+\dfrac{y^2}{8}=1$.
}
\end{ex}

\begin{ex}%[0H9H5-4]%[Dự án đề kiểm tra Toán 10 GHKI NH23-24- Lam Nguyen]%[THPT Le Quý Đôn - TPHCM]
Cho đường hypebol có phương trình $(H): 9x^2-3y^2=1$. Khoảng cách giữa hai tiêu điểm của hypebol là
\choice
{$0$}
{$\dfrac{\sqrt{2}}{3}$}
{\True $\dfrac{4}{3}$}
{$\dfrac{2}{3}$}
\loigiai{
Phương trình chính tắc của hypebol $(H): \dfrac{x^2}{\tfrac{1}{9}}-\dfrac{y^2}{\tfrac{1}{3}}=1$.\\
Vậy ta có $a^2=\dfrac{1}{9}$, $b^2=\dfrac{1}{3}$.\\
$c^2=b^2+a^2=\dfrac{1}{3}+\dfrac{1}{9}=\dfrac{4}{9} \Rightarrow c=\dfrac{2}{3}$.\\
Vậy khoảng cách giữa hai tiêu điểm của hypebol là $2c=\dfrac{4}{3}$.

}
\end{ex}
\TNTF
\begin{ex}%[0D8H2-4]
Một câu lạc bộ cờ vua có $10$ bạn nam và $7$ bạn nữ. Huấn luyện viên muốn chọn $4$ bạn đi thi đấu cờ vua.
\choiceTF
{\True Có $210$ cách chọn $4$ bạn nam}
{Có $2308$ cách chọn $4$ bạn không phân biệt nam, nữ}
{\True Có $945$ cách chọn $4$ bạn, trong đó có $2$ bạn nam và $2$ bạn nữ}
{\True Có $2345$ cách chọn $4$ bạn, trong đó có nhiều nhất $3$ bạn nữ}
\loigiai{
\begin{itemchoice}
\itemch Đúng vì \\
Chọn $4$ bạn nam trong $10$ bạn nam là $\mathrm{C}_{10}^4=210$ cách.
\itemch Sai vì \\
Chọn $4$ bạn không phân biệt nam nữ từ $17$ bạn là tổ hợp chấp $4$ của $17$ phần tử, nên số cách chọn là $\mathrm{C}_{17}^{4}=2380$ cách.
\itemch Đúng vì\\
Chọn $2$ bạn nam trong $10$ nam, có $\mathrm{C}_{10}^2=45$ cách.
Chọn $2$ bạn nữ trong $7$ nữ, có $\mathrm{C}_7^2=21$ cách.\\
Vậy số cách chọn $4$ bạn, có $2$ nam, $2$ nữ là $45\cdot 21=945$ cách.
\itemch Đúng vì \\
Chọn $4$ nữ từ $7$ là $\mathrm{C}_7^4$ cách. \\
Số cách chọn tùy ý là $\mathrm{C}_{17}^4=2380$ cách. \\
Suy ra số cách thỏa đề là $2345$.
\end{itemchoice}
}
\end{ex}

\begin{ex}%[Dự án 18 - TeamTeXHoa - Phạm Quốc Toàn]%[0H9H4-5]
Trong mặt phẳng tọa độ $Oxy$, cho hai điểm $A \left(1;-4 \right)$, $B \left(-2;0 \right)$ và đường thẳng \\
$d \colon 2x-4y+1=0$. Các mệnh đề sau đúng hay sai?
\choiceTF
{Điểm $A$, $B$ cách đều đường thẳng $d$}
{Tọa độ tâm của đường tròn đi qua hai điểm $A$, $B$ và có tâm thuộc đường thẳng $d$ là $I \left(1;\dfrac{3}{4}\right)$}
{\True Phương trình đường tròn $\left(C \right)$ đi qua $A$, $B$ và có tâm thuộc đường thẳng $d$ là \\
$4x^2+4y^2-60x-32y-136=0$}
{\True Giá trị nhỏ nhất của $OM$ với $M$ là điểm chuyển động trên đường tròn là $\dfrac{5\sqrt{17}-17}{2}$}
\loigiai{
\begin{itemchoice}
\itemch Khoảng cách từ điểm $A$ đến đường thẳng $d$ là $d\left(A, d\right)=\dfrac{\left| 2\cdot 1-4 \cdot \left(-4 \right)+1\right|}{\sqrt{2^2+ \left(-4 \right)^2}}=\dfrac{19\sqrt{5}}{10}$. \\
Khoảng cách từ điểm $B$ đến đường thẳng $d$ là $d\left(B, d\right)=\dfrac{\left| 2\cdot \left(-2 \right)-4 \cdot 0 +1\right|}{\sqrt{2^2+ \left(-4 \right)^2}}=\dfrac{3\sqrt{5}}{10}$. \\
Vậy hai điểm $A$, $B$ không cách đều đường thẳng $d$.
\itemch Đường tròn có tâm thuộc đường thẳng $d \colon 2x-4y+1=0$ nên gọi $I$ là tâm đường tròn thì $I\left(2a-\dfrac{1}{2};a\right)$.\\
Mặt khác đường tròn đi qua $A$, $B$ nên $IA=IB\Leftrightarrow \left| \overrightarrow{IA}\right|=\left| \overrightarrow{IB}\right|$ \\
$\Leftrightarrow \sqrt{\left(1-2a+\dfrac{1}{2}\right)^2+ \left(-4-a \right)^2}=\sqrt{\left(-2-2a+\dfrac{1}{2}\right)^2+ \left(0-a \right)^2}\Leftrightarrow a=4$. \\
Vậy đường tròn đi qua $A$, $B$ và có tâm thuộc đường thẳng $d$ có tọa độ tâm là $I\left(\dfrac{15}{2};4\right)$.
\itemch $R=IB=\left| \overrightarrow{IB}\right|=\sqrt{\left(-2-\dfrac{15}{2}\right)^2+ \left(0-4 \right)^2}=\dfrac{5\sqrt{17}}{2}$. \\
Phương trình đường tròn đi qua $A$, $B$ và có tâm thuộc đường thẳng $d$ là \\
$\left(x-\dfrac{15}{2}\right)^2+ \left(y-4 \right)^2=\left(\dfrac{5\sqrt{17}}{2}\right)^2 \Leftrightarrow 4x^2+4y^2-60x-32y-136=0$.
\itemch Gọi $G$, $K$ lần lượt là giao điểm của đường thẳng $IO$ và đường tròn.
\begin{center}
\begin{tikzpicture}[>=stealth, x=1.0 cm, y=1.0 cm]
\draw[->] (-3,0) --(0,0) node[below right]{$O$} -- (4,0) node[below]{$x$};
\draw[->](0,-3)--(0, 4.5) node[right]{$y$};
\draw(1,1) circle (3); %(0,0) là tọa độ tâm, 1 là bán kính
\draw(1,1) -- (-1.12, -1.12) (1, 1) -- (3.11, 3.13) (0, 0) -- (1, 4) -- (1, 1);
\draw (1,1) circle (1pt) node[below]{$I$};
\draw (1,4) circle (1pt) node[above]{$M$};
\draw (-1.12, -1.12) circle (1pt) node[below left]{$G$};
\draw (3.11, 3.13) circle (1pt) node[above right]{$K$};
\end{tikzpicture}
\end{center}
Với ba điểm $I$, $O$, $M$ ta có \\
$IM-IO \le OM \le IM+IO \Rightarrow IG-IO \le OM \le IK+IO \Rightarrow OG \le OM \le OK$. \\
Trong đó $OG=IG-IO=\dfrac{5\sqrt{17}}{2}-\dfrac{17}{2}=\dfrac{5\sqrt{17}-17}{2}$. \\
Vậy $OM$ đạt giá trị nhỏ nhất là $\dfrac{5\sqrt{17}-17}{2}$ khi $M$ trùng $G$.
\end{itemchoice}
}
\end{ex}

\TNSA
\begin{ex}%[0D8H3-2]%[Dự án tex đề CKII - Đợt 15]%[Đình Trí]
Tìm hệ số chứa $ x^4 $ trong khai triển $ (2x-1)^5 $.
\loigiai
{
Ta có
\begin{eqnarray*}
(2x-1)^5 &=&\mathrm{C}_5^0(2x)^5-\mathrm{C}_5^1(2x)^4+\mathrm{C}_5^2(2x)^3-\mathrm{C}_5^3(2x)^2+\mathrm{C}_5^4(2x)-\mathrm{C}_5^5\\
&=& 32x^5-80x^4+80x^3-40x^2+10x-1.
\end{eqnarray*}
Dựa vào khai triển suy ra hệ số chứa $ x^4 $ là $ -80 $.
}
\end{ex}

\begin{ex}%[0D8H2-5]%[Dự án Toán thực tế]%[Vũ Hồng Toàn]
Trong trò chơi \lq\lq Chiếc nón kì diệu\rq\rq\ chiếc kim của bánh xe có thể dừng lại ở một trong $7$ vị trí với các khả năng như nhau. Hỏi có bao nhiêu cách để trong ba lần quay, chiếc kim của bánh xe đó lần lượt dừng lại ở ba vị trí khác nhau?
\begin{center}
\begin{tikzpicture}[declare function={r=2;n=7;goc=360/n;}]
\draw (0,0) coordinate (O) circle (r);
\foreach \i in {1,...,7}{\draw (O)--({\i*goc}:r);
\node[circle,fill=cyan!30, inner sep=.4,text=blue] at ({\i*goc + goc/2}:0.75*r) {$\i$};
}
\draw[fill=red!50]
(30:0.5*r)--($ (O) + (80:{0.1*r}) $) arc[start angle=90, end angle=300, radius=0.1*r]-- cycle;
\end{tikzpicture}
\end{center}
\shortans{$210$}
\loigiai{
Lần quay thứ nhất, bánh xe có $7$ cách dừng kim.\\
Lần quay thứ hai, bánh xe có $6$ cách dừng kim (khác vị trí của lần thứ nhất).\\
Lần quay thứ ba, bánh xe có $5$ cách dừng kim (khác vị trí của lần thứ nhất và lần thứ hai).\\
Vậy có $7\cdot 6\cdot 5=210$ cách.
}
\end{ex}

\begin{ex}%[0H9V4-1]
Cho $A(-1;0)$, $B(2;4)$ và $C(4;1)$. Biết rằng tập hợp các điểm $M$ thoả mãn $3MA^2+MB^2=2MC^2$ là một đường tròn $(C)$. Bán kính của $(C)$ là $R=\dfrac{\sqrt{m}}{n}$ với $m$, $n$ là các số nguyên tố. Khi đó $m+n$ bằng bao nhiêu?
\shortans[oly]{109}
\loigiai{
Gọi $M(x;y)$, ta có
\begin{eqnarray*}
&&3MA^2+MB^2=2MC^2\\
&\Leftrightarrow& 3\left[(x+1)^2+y^2\right]+(x-2)^2+(y-4)^2=2\left[(x-4)^2+(y-1)^2\right]\\
&\Leftrightarrow& 3x^2+3y^2+6x+3+x^2+y^2-4x-8y+20=2x^2+2y^2-16x-4y+34\\
&\Leftrightarrow& 2x^2+2y^2+18x-4y-11=0\\
&\Leftrightarrow& x^2+y^2+9x-2y-\dfrac{11}{2}=0. \quad (*)
\end{eqnarray*}
Đặt $a=-\dfrac{9}{2}$, $b=1$, $c=-\dfrac{11}{2}$.\\
Ta có $a^2+b^2-c=\dfrac{107}{4} > 0$ nên $\left(*\right)$ là phương trình của một đường tròn $(C)$.\\
Bán kính của $(C)$ là $R=\dfrac{\sqrt{107}}{2}=\dfrac{\sqrt{m}}{n}$.\\
Khi đó $m+n=107+2=109$
}
\end{ex}
\begin{ex}%[0H9V5-0]
\immini
{
Một công viên có dạng hình elip có độ dài trục lớn bằng $100$ m, độ dài trục bé bằng $80$ m. Người ta muốn trồng hoa trong một hình chữ nhật nội tiếp elip, phần còn lại sẽ trồng cỏ (như hình vẽ). Diện tích trồng hoa lớn nhất bằng bao nhiêu mét vuông?
}
{
\begin{tikzpicture}[line join = round, line cap = round,>=stealth,scale=.8,font=\footnotesize]
\def\ra{3cm}
\def\rb{2cm}
\path
(0,0) coordinate (A')++(0:2*\ra)coordinate(A)
($(A')!.5!(A)$)coordinate(O)
($(O)+(90:\rb)$)coordinate(B)++(-90:2*\rb)coordinate(B')
($(O)+(50:3cm and 2cm)$) coordinate (M)
($(O)+(130:3cm and 2cm)$) coordinate (N)
($(O)+(-130:3cm and 2cm)$) coordinate (P)
($(O)+(-50:3cm and 2cm)$) coordinate (Q)
;
\draw
(M)--(N)--(P)--(Q)--cycle
(A')--(A) (B')--(B)
;
\fill[pattern=dots,pattern  color=gray]
(M)--(N)--(P)--(Q)--cycle
;
\draw (O) ellipse (3cm and 2cm);
\foreach \point/\angle in {A'/190,A/-10,B/80,B'/-100}{
\draw[fill=black](\point)circle(1.2pt)node[shift={(\angle:2.5mm)}]{$\point$};
}
\foreach \point in {M,N,P,Q,O}{
\draw[fill=black](\point)circle(1.2pt);
}
\end{tikzpicture}
}
\shortans{$4000$}
\loigiai{
Phương trình chính tắc của elip $(E)\colon \dfrac{x^2}{a^2}+\dfrac{y^2}{b^2}=1, (a > b > 0)$.\\
Theo giả thiết ta có $\heva{&
2a=100\Rightarrow a=50\\&
2b=80\Rightarrow b=40.}$\\
Suy ra phương trình elip $(E)\colon \dfrac{x^2}{2500}+\dfrac{y^2}{1600}=1$.\\
Gọi $M(x_M;y_M)$ là đỉnh hình chữ nhật với $x_M> 0$, $y_M> 0$.\\
Ta có $\dfrac{x_M^2}{2500}+\dfrac{y_M^2}{1600}=1$.\\
Diện tích hình chữ nhật là
\[S=2x_M\cdot 2y_M=4x_My_M=4000\cdot 2\cdot \dfrac{x_M}{50}\cdot \dfrac{y_M}{40}\le 4000\left(\dfrac{x_M^2}{2500}+\dfrac{y_M^2}{1600}\right)=4000\  \mathrm{(m^2)}.\]
}
\end{ex}

\TL
\begin{ex}%[Phan Quốc Trí]%[0D8H3-4]
Tìm hệ số của $x^3$ trong khai triển biểu thức $P(x)=x(1-x)^4$.
% \shortans[]{$6$}
\loigiai{
Ta có $P(x)=x(1-x)^4=x\left(\mathrm{C}_4^0 +\mathrm{C}_4^1 (-x)+\mathrm{C}_4^2 (-x)^2+\mathrm{C}_4^3 (-x)^3+\mathrm{C}_4^4 (-x)^4\right)$.\\
Suy ra số hạng chứa $x^3$ là $\mathrm{C}_4^2 x^3 =6x^3$.\\
Do đó hệ số của $x^3$ trong khai triển biểu thức $P(x)=x(1-x)^4$ thành đa thức bằng $6$.
}
\end{ex}
\begin{ex}%[0D8V2-4]
Một cuộc họp có sự tham gia của $6$ nhà Toán học trong đó có 4 nam và $2$ nữ, $7$ nhà Vật lý trong đó có $3$ nam và $4$ nữ và $8$ nhà Hóa học trong đó có $4$ nam và $4$ nữ. Có bao nhiêu cách lập một ban thư kí gồm $4$ nhà khoa học, trong đó phải có đủ cả $3$ lĩnh vực và có cả nam lẫn nữ.
\loigiai{
\begin{itemize}
\item Ta tìm số cách chọn ra được $4$ nhà khoa học có đầy đủ cả $3$ lĩnh vực.
\begin{itemize}
\item Số cách chọn 2 nhà Toán học, 1 nhà Vật lý, 1 nhà Hóa học là $\mathrm{C}_6^2\cdot\mathrm{C}_7^1\cdot\mathrm{C}_8^1=840$.
\item Số cách chọn 1 nhà Toán học, 2 nhà Vật lý, 1 nhà Hóa học là $\mathrm{C}_6^1\cdot\mathrm{C}_7^2\cdot\mathrm{C}_8^1=1008$.
\item Số cách chọn 1 nhà Toán học, 1 nhà Vật lý, 2 nhà Hóa học là $\mathrm{C}_6^1\cdot\mathrm{C}_7^1\cdot\mathrm{C}_8^2=1176$.
\end{itemize}
Suy ra có $840+1008+1176=3024$ cách chọn ra được $4$ nhà khoa học có đầy đủ cả $3$ lĩnh vực.
\item Ta tìm số cách chọn $4$ nhà khoa học đủ cả $3$ lĩnh vực mà trong đó chỉ có nam hoặc chỉ có nữ.
\begin{itemize}
\item Số cách chọn chỉ có nam là $\mathrm{C}_4^2\cdot\mathrm{C}_3^1\cdot\mathrm{C}_4^1+\mathrm{C}_4^1\cdot\mathrm{C}_3^2\cdot\mathrm{C}_4^1+\mathrm{C}_4^1\cdot\mathrm{C}_3^1\cdot\mathrm{C}_4^2=192$.
\item Số cách chọn chỉ có nữ là $\mathrm{C}_2^2\cdot\mathrm{C}_4^1\cdot\mathrm{C}_4^1+\mathrm{C}_2^1\cdot\mathrm{C}_4^2\cdot\mathrm{C}_4^1+\mathrm{C}_2^1\cdot\mathrm{C}_4^1\cdot\mathrm{C}_4^2=112$.\\
Suy ra có $192+112=304$ cách chọn $4$ nhà khoa học đủ cả $3$ lĩnh vực mà trong đó chỉ có nam hoặc chỉ có nữ.
\end{itemize}
\end{itemize}
Vậy số cách chọn ra được $4$ nhà khoa học có đầy đủ cả $3$ lĩnh vực, trong đó có cả nam lẫn nữ là $3024-304=2720$ cách.
}
\end{ex}

\begin{ex}%[10-11-12EX-HK1-2425]%[Huỳnh Xuân Tín]%[0H9V4-2]
Trong mặt phẳng $Oxy$, cho điểm $I(-2; 3)$ và đường thẳng $d\colon 3x+4y-1=0$.
\begin{enumerate}
\item Tính khoảng cách từ điểm $I$ đến đường thẳng $d$.
\item Viết phương trình đường tròn tâm $I$ và tiếp xúc với đường thẳng $d$.
\end{enumerate}
\loigiai{
\begin{enumerate}
\item Khoảng cách từ điểm $I$ đến đường thẳng $d$ là
\[\mathrm{d}(I,d)=\dfrac{|3\cdot(-2)+4\cdot3-1|}{\sqrt{3^2+4^2}}=1.\]
\item Đường tròn tâm $I$ và tiếp xúc với đường thẳng $d$ có bán kính là $1$ nên có phương trình là
\[(x+2)^2+(y-3)^2=1.\]
\end{enumerate}
}
\end{ex}

\begin{ex}%[0D0C2-7]%[Dự án đề kiểm tra Toán 11 GHKII NH23-24- Tổng Nguyễn]%[THPT Đống Đa - Hà Nội]
Ba bạn Bình, Duy, Nam mỗi bạn viết ngẫu nhiên một số tự nhiên thuộc đoạn $[1;16]$ được kí hiệu theo thứ tự $a,b,c$ rồi lập phương trình bậc hai $ax^2+2bx+c=0$. Có bao nhiêu bộ số $(a;b;c)$ sao cho phương trình lập được có nghiệm kép.
% \shortans[]{$129$}
\loigiai{ Phương trình $ax^2+2bx+c=0$ có nghiệm kép khi  $\Delta' =b^2-ac=0\Leftrightarrow b^2=ac$.\\
Vì $a$, $b$, $c$ là các số nguyên thuộc đoạn $[1;16]$  thỏa mãn $b^2=ac$ nên các bộ $(b,ac)$ là\\
$(1,1)$, $(2,4)$, $(3,9)$, $(4,16)$, $(5,25)$, $(6,36)$, $(7,49)$, $(8,64)$, $(9,81)$, $(10,100)$, $(11,121)$, $(12,144)$, $(13,169)$, $(14,196)$, $(15,225)$, $(16,216)$.
\begin{itemize}
\item Vì $1^2=1\cdot1$ nên bộ $(1,1)$ cho $1$ bộ $(a,b,c)$ là $(1,1,1)$.
\item  Vì $2^2=4=2\cdot2=4\cdot1=1\cdot 4$ nên bộ $(2,4)$ cho $3$ bộ $(a,b,c)$ thỏa $b^2=ac$.
\item  Vì $3^2=9=3\cdot3=9\cdot1=1\cdot 9$ nên bộ $(3,9)$ cho $3$ bộ $(a,b,c)$ thỏa $b^2=ac$.
\item  Vì $4^2=16=4\cdot4=16\cdot1=1\cdot 16=2\cdot8=8\cdot2$ nên bộ $(4,16)$ cho $5$ bộ $(a,b,c)$ thỏa $b^2=ac$.
\item  Vì $5^2=25=5\cdot5$ nên bộ $(5,25)$ cho $1$ bộ $(a,b,c)$ thỏa $b^2=ac$.
\item  Vì $6^2=36=6\cdot6=12\cdot3=3\cdot12=9\cdot4=4\cdot9$ nên bộ $(6,36)$ cho $5$ bộ $(a,b,c)$ thỏa $b^2=ac$.
\item  Vì $7^2=49=7\cdot7$ nên bộ $(7,49)$ cho $1$ bộ $(a,b,c)$ thỏa $b^2=ac$.
\item  Vì $8^2=64=8\cdot8=4\cdot16=16\cdot4$ nên bộ $(8,64)$ cho $3$ bộ $(a,b,c)$ thỏa $b^2=ac$.
\item  Vì $9^2=81=9\cdot9$ nên bộ $(9,81)$ cho $1$ bộ $(a,b,c)$ thỏa $b^2=ac$.
\item  Vì $10^2=100=10\cdot10$ nên bộ $(10,100)$ cho $1$ bộ $(a,b,c)$ thỏa $b^2=ac$.
\item  Vì $11^2=121=11\cdot11$ nên bộ $(11,121)$ cho $1$ bộ $(a,b,c)$ thỏa $b^2=ac$.
\item  Vì $12^2=144=12\cdot12=9\cdot16=16\cdot9$ nên bộ $(12,144)$ cho $3$ bộ $(a,b,c)$ thỏa $b^2=ac$.
\item Mỗi bộ $(b,ac)$ với $13\le b\le 16$ cho đúng $1$ bộ $(a,b,c)$ thỏa $b^2=ac$.
\end{itemize}
Như vậy có tất cả $1+3+3+5+1+5+1+3+1+1+1+3+4\cdot1=32$.\\
% Không gian mẫu của phép thử trên có số phần tử là $n(\Omega)=16\cdot 16\cdot 16=4096$.\\
% Xác suất cần tìm là $P=\dfrac{32}{4096}=\dfrac{1}{128}$. Suy ra $m=1,n=128$.\\
% Vậy $m+n=1+128=129$.

}
\end{ex}

\Closesolutionfile{ans}
\Closesolutionfile{ansbook}
